\honest{Detached Lever Fallacy}{Eliezer Yudkowsky}{2008}

I name this fallacy after an old science-fiction show that I never saw; someone at a convention told me about it. The heroes capture an alien ship that can pass through asteroids. The captain finds a lever on the bridge, pries it off, and takes it back to his own ship, which can then pass through asteroids too. \nb{The show could not be traced. A commenter's suggestion, a 1994 episode of \textsc{Star Trek: The Next Generation}, matches only loosely.}

People still program AI with semantic networks such as ``(apple is-a fruit).'' You have seen, touched, bought, sliced and tasted apples. Your temporal cortex somehow stores a generalized image that recognizes new apples from new angles, and your motor cortex and cerebellum store programs for using them. Writing ``apple,'' five ASCII characters, pulls a lever on another person's similar machinery. Written into a program, it is ``just a lever.'' \nb{Drew McDermott mocked such networks in a 1976 paper, which the post cites and which ``Truly Part of You'' introduced. Stevan Harnad later called this the symbol grounding problem: meaning ``parasitic on the meanings in our heads.''} Silicon could have the machinery someday, but ``you can't just pry the control lever off the bridge!'' The fallacy tempts us because the lever is visible and varies, while the machinery is hidden and always there. Most people do not know their temporal cortex exists. Since perception is largely the perception of differences, machinery that never switches off is not noticed as a requirement.

This is why you cannot make an AI kind by giving it loving parents, ``As I've often heard proposed.'' \nb{The post names no one. Turing proposed in 1950 to educate a ``child machine'' with rewards and punishments, and MIT's Cog robot was built to learn from people ``as human infants do.''} In biology, growing a fur coat in response to cold takes more genetic machinery than growing one regardless, since the cold must be sensed. \nb{The literature on plasticity agrees.} Forgetting this leads to ``Lamarckian delusions,'' and there were ``various slap-fights'' of this sort. \nb{None is named. Trofim Lysenko's claims about cold-treated wheat are a close match.}

Culture, too, is absorbed through inherited responses; try raising a fish as a Mormon or sending a lizard to college. Children who hear only a mixture of pseudo-languages grow up speaking a language with syntax. \nb{This is Bickerton's creole hypothesis, which is disputed. Nicaraguan Sign Language, created by deaf children, fits the post's point better.} Soviet posters could not raise selfless workers. No known childhood environment makes humans selfless. Game theory explains an innate response of returning kindness for kindness and hatred for hatred, provided the kindness does not look too unconditional, hence spoiled children and the testing of limits. Abused children ``have a much higher probability'' of abusing their own, though many break the loop. For more, I point to Tooby and Cosmides and to Pinker's \textsc{The Blank Slate}. \nb{The evidence is mixed. A 1987 review found that the unqualified belief is unfounded. In a 2015 prospective study, adults abused as children did not report abusing their own children more than others did, though their children reported more neglect and sexual abuse.}

Raise a baby AI with kind but strict parents and you pull levers on machinery it does not have, built in humans by millions of years of selection. We absorb culture faithfully only because we are humans absorbing a human culture, and humans raised in an alien culture would probably make it more human. The human version is sloppy and works only under particular conditions. It would be stupid and dangerous to build a naughty AI that tests its boundaries and has to be spanked, rather than just asking.

Will programmers write, line by line, the code by which an AI that feels slighted conceives a hatred of its makers, a sullen-teenager AI? Unconditional niceness is easier to program than niceness conditional on upbringing, so if you cannot do the first, you cannot do the second. A paperclip maximizer raised with love stays a paperclip maximizer. Kindness is not sneezed into an AI by contagion. Getting information from the environment is possible, but structuring the response so the AI ends up where you want is the major problem. The word ``learning'' sounds as if the magic is in the environment, when the magic is in the structured response. ``The blank slate is a chimera.'' \nb{Turing had described the child brain as ``Rather little mechanism, and lots of blank sheets,'' but also said the child program must be found by experiment, which agrees with the post.}

The world is deeper than it appears: every word in print and everything we teach children are surface levers on hidden machinery, and because they are all that varies, they seem to be all that exists. People near AI ``usually'' build imitation levers and are surprised when nothing happens. \nb{No case is given. Robin Hanson objected in the comments that this slandered a generation of researchers; the author replied that ``a very substantial fraction'' made the mistake.} So when someone proposes to raise an AI in a loving family, or among liberal democratic values, think of a lever pried off the bridge.
